% (User input window)

\paragraph{\texorpdfstring{$\delta_q$}{deltaq} 的几何化：谱维数量纲、二次通道极大指数与样本阈值}
沿用注 \ref{rem:pom-collision-kernel-rh-phase} 的记号，\(\delta_q=\log(\Lambda_q^2/r_q)\) 描述非 Perron 通道的二次指数尺度 \(\Lambda_q^{2m}\) 相对 bulk 尺度 \(r_q^{m}\) 的竞争强度。
本段把该越线指数在 \(\varphi\)-标度下几何化，并给出其作为“二次观测最强相对指数”与“最小样本指数阈值”的可审计表述。

\begin{definition}[bulk/bdry 谱维数与相关余维]\label{def:pom-deltaq-spectral-dimensions}
对每个整数 \(q\ge 2\)，令 \(r_q\) 与 \(\Lambda_q\) 分别为 \(A_q\) 的 Perron 根与非 Perron 最大模根（注 \ref{rem:pom-collision-kernel-rh-phase}）。
在几何标度 \(\varepsilon_m=\varphi^{-m}\)（命题 \ref{prop:pom-renyi-dimension-spectrum}）下，定义
\[
d_{\mathrm{bulk}}(q):=\frac{\log r_q}{\log\varphi},\qquad
d_{\mathrm{bdry}}(q):=\frac{\log \Lambda_q}{\log\varphi}.
\]
并定义“相关边界余维”
\[
\mathrm{codim}_{\mathrm{corr}}(q):=\frac{\log r_q-2\log\Lambda_q}{\log\varphi}.
\]
\end{definition}

\begin{proposition}[\texorpdfstring{$\delta_q$}{deltaq} 的几何恒等式]\label{prop:pom-deltaq-spectral-dimension-identity}
在定义 \ref{def:pom-deltaq-spectral-dimensions} 的记号下，成立严格恒等式
\[
\boxed{\;
\frac{\delta_q}{\log\varphi}
=2d_{\mathrm{bdry}}(q)-d_{\mathrm{bulk}}(q)
\;}
\qquad\text{以及}\qquad
\boxed{\;
\mathrm{codim}_{\mathrm{corr}}(q)=-\frac{\delta_q}{\log\varphi}\; }.
\]
\end{proposition}
\begin{proof}
由 \(\delta_q=2\log\Lambda_q-\log r_q\)，两边同除以 \(\log\varphi\) 即得第一式。
第二式由 \(\mathrm{codim}_{\mathrm{corr}}\) 的定义与第一式直接整理得到。证毕。
\end{proof}

\begin{corollary}[半阈值判据与越线相]\label{cor:pom-deltaq-half-threshold-criterion}
在同一记号下，下列条件严格等价：
\begin{enumerate}
  \item \(\delta_q=0\)；
  \item \(d_{\mathrm{bdry}}(q)=\tfrac12 d_{\mathrm{bulk}}(q)\)；
  \item \(\beta^{\mathrm{RH}}_q=\tfrac12\)（注 \ref{rem:pom-collision-kernel-rh-phase}）。
\end{enumerate}
并且 \(\delta_q>0\)（分别 \(\delta_q<0\)）当且仅当
\(
d_{\mathrm{bdry}}(q)>\tfrac12 d_{\mathrm{bulk}}(q)
\)
（分别 \(<\)）。
\end{corollary}
\begin{proof}
由命题 \ref{prop:pom-deltaq-spectral-dimension-identity} 与 \(\beta^{\mathrm{RH}}_q=\log\Lambda_q/\log r_q\) 直接代入整理即可。证毕。
\end{proof}

\input{sections/body/pom/parts/part__pom-spectral-condensation-index-gLambda}

\begin{proposition}[中心化二次观测的极大相对指数]\label{prop:pom-deltaq-quadratic-observable-max-exponent}
设 \(A\) 为实方阵，谱半径为 \(r>0\)，并令
\[
\Lambda:=\max_{\mu\in\sigma(A)\setminus\{r\}}|\mu|.
\]
假设存在 Perron--Frobenius 型分解
\[
A^m=r^m\Pi+R_m,\qquad \|R_m\|=O(\Lambda^m),
\]
其中 \(\Pi\) 为 Perron 投影。
对任意线性泛函 \(\mathcal L(\cdot)\) 定义线性观测 \(Z_m:=\mathcal L(A^m)\)，并令 \(c:=\mathcal L(\Pi)\) 与中心化观测 \(\widetilde Z_m:=Z_m-c\,r^m\)。
则：
\begin{enumerate}
  \item 对任意 \(\mathcal L\) 有 \(|\widetilde Z_m|=O(\Lambda^m)\)，从而
  \[
  \limsup_{m\to\infty}\frac{1}{m}\log|\widetilde Z_m|\le \log\Lambda.
  \]
  \item 若 \(\mathcal L\) 在达到模 \(\Lambda\) 的谱投影上不湮灭，则存在常数 \(c_\ast>0\) 使沿无穷多 \(m\) 有 \(|\widetilde Z_m|\ge c_\ast \Lambda^m\)，从而
  \[
  \limsup_{m\to\infty}\frac{1}{m}\log|\widetilde Z_m|=\log\Lambda.
  \]
  \item 对任意两组线性泛函 \(\mathcal L_1,\mathcal L_2\) 与相应中心化观测 \(\widetilde Z^{(i)}_m\)，其二次/双线性观测
  \[
  Q_m:=(\widetilde Z_m)^2\quad\text{或}\quad Q_m:=\widetilde Z^{(1)}_m\,\widetilde Z^{(2)}_m
  \]
  满足 \(|Q_m|=O(\Lambda^{2m})\)，并在两泛函均不退化时有
  \[
  \limsup_{m\to\infty}\frac{1}{m}\log|Q_m|=2\log\Lambda.
  \]
  因此相对主尺度 \(r^m\) 的最大指数为
  \[
  \boxed{\;
  \limsup_{m\to\infty}\frac{1}{m}\log\frac{|Q_m|}{r^m}
  =\log\!\left(\frac{\Lambda^2}{r}\right)\; }.
  \]
\end{enumerate}
\end{proposition}
\begin{proof}
将分解 \(A^m=r^m\Pi+R_m\) 代入 \(Z_m=\mathcal L(A^m)\) 得
\(
\widetilde Z_m=\mathcal L(R_m)
\)，从而 (1) 由 \(\|R_m\|=O(\Lambda^m)\) 立得。
(2) 可由谱投影展开（或 Jordan 分解）给出：若 \(\mathcal L\) 在模 \(\Lambda\) 的谱分量上不为零，则该分量沿一无穷子列不能被抵消，从而 \(\limsup\) 达到 \(\log\Lambda\)。
(3) 由 (1) 立即得到 \(Q_m=O(\Lambda^{2m})\)，而在非退化时沿同一子列取得下界，故 \(\limsup\) 为 \(2\log\Lambda\)；最后一式只是把 \(2\log\Lambda-\log r\) 改写为 \(\log(\Lambda^2/r)\)。证毕。
\end{proof}

\begin{corollary}[越线相的最小样本指数阈值]\label{cor:pom-deltaq-sample-complexity-threshold}
沿命题 \ref{prop:pom-deltaq-quadratic-observable-max-exponent} 的记号，设存在可采样随机观测 \(X_m\) 满足
\[
\mathbb E[X_m]=c\,r^{m/2},\qquad \mathrm{Var}(X_m)=\Theta(\Lambda^{2m})
\]
（其中 \(\Lambda^{2m}\) 表示中心化二次通道的指数尺度）。
令 \(\widehat X_{m,N}\) 为 \(N\) 次独立采样的样本均值估计，则相对均方误差满足
\[
\frac{\mathbb E[(\widehat X_{m,N}-\mathbb E[X_m])^2]}{\mathbb E[X_m]^2}
=\frac{\mathrm{Var}(X_m)}{N\,\mathbb E[X_m]^2}
\asymp
\frac{1}{N}\exp\!\left(m\log\!\left(\frac{\Lambda^2}{r}\right)\right).
\]
因此要使相对误差在 \(m\to\infty\) 下保持 \(o(1)\)，必要且（在常数意义下）充分的样本量标度为
\[
N\gtrsim \exp\!\left(m\log\!\left(\frac{\Lambda^2}{r}\right)\right).
\]
应用到 \(A_q\) 即得：\(\delta_q=\log(\Lambda_q^2/r_q)\) 给出从“多项式可控”到“指数门槛”的样本复杂度相变指数。
\end{corollary}
\begin{proof}
由独立样本均值的方差公式
\(
\mathrm{Var}(\widehat X_{m,N})=\mathrm{Var}(X_m)/N
\)
直接代入并整理得到。证毕。
\end{proof}

\begin{remark}[最小边界可见阶与离散证书化]\label{rem:pom-deltaq-minimal-visible-order-auditability}
由表 \ref{tab:fold_collision_kernel_rh_scan_q2_17} 的符号变化
\(
\delta_2<0,\ \delta_3<0,\ \delta_4>0
\)
可把 \(q=4\) 识别为复制阶方向上的最小越线点（亦即最小的“边界主导阶”，注 \ref{rem:pom-collision-kernel-rh-phase}）。
\par
同时，\((r_q,\Lambda_q,\delta_q)\) 可在原则上完全脱离任何外部拟合而由有限数据审计：由 Hankel 秩定理（定理 \ref{thm:hankel-rank-minimal}）与最小实现函子 \(\mathsf{MinH}\)（注 \ref{rem:pom-minh-functor}），序列 \(S_q(m)\) 的最小递推多项式 \(P_q(\lambda)\) 在同构意义下由有限 Hankel 数据规范确定；而 \(r_q\) 与 \(\Lambda_q\) 可直接由 \(P_q\) 的根模结构抽取，从而 \(\delta_q=\log(\Lambda_q^2/r_q)\) 获得离散可复核的谱证书（参见算法 \ref{alg:pom-recurrence-polynomial-spectral-extraction}）。
\end{remark}

